\documentclass[]{article}

\usepackage{amsmath}
\usepackage{multirow}
\usepackage{natbib}
\usepackage{graphicx} 


\begin{document}

\subsection{Initial conditions}
The simulation follows the evolution of $512^3$ particles in a periodic cosmological box. The initial conditions were generated at a starting redshift of $z=127$ using the Zel'dovich Approximation. We first computed the linear matter power spectrum, $P(k)$, at this redshift using the Code for Anisotropies in the Microwave Background (CAMB), assuming the fiducial cosmology from the Quijote simulations ($\Omega_m=0.3175$, $h=0.6711$, $n_s=0.9624$, $\sigma_8=0.834$). A Gaussian random field was then realized in Fourier space, which was inverse-transformed to generate the displacement and velocity fields. Particles were displaced from a uniform grid according to these fields, setting up the initial state for the simulation.

\subsection{GPU-accelerated Particle-Mesh simulation}
We developed a Particle-Mesh (PM) gravity solver using the NVIDIA Warp framework, a Python-based toolkit for writing high-performance GPU kernels. The simulation loop consists of three main stages: mass assignment, solving for the gravitational potential, and updating particle positions and velocities.

\subsubsection{Mass assignment and force interpolation}
The particle mass is assigned to a regular $512^3$ grid using a Cloud-in-Cell (CIC) scheme. This step was implemented as a custom Warp kernel that iterates over all particles in parallel. To handle multiple particles contributing to the same grid cell concurrently, we used atomic add operations (`wp.atomic\_add`) to ensure thread-safe updates to the density field. After solving for the potential, the gravitational force at each particle's position is calculated by interpolating from the grid using the same CIC weighting scheme in a gather-style operation, also implemented as a custom Warp kernel.

\subsubsection{Poisson solver}
The gravitational potential, $\Phi$, is computed by solving the Poisson equation in comoving coordinates:
\begin{equation}
    \nabla^2 \Phi = \frac{3}{2} \Omega_m H_0^2 \frac{\delta}{a}
\end{equation}
where $\delta = (\rho - \bar{\rho}) / \bar{\rho}$ is the matter density contrast, $a$ is the scale factor, $\Omega_m$ is the matter density parameter, and $H_0$ is the Hubble constant. This equation is solved efficiently in Fourier space. The density grid is transformed using a 3D real-to-complex Fast Fourier Transform (FFT). In Fourier space, the equation becomes algebraic, and the potential $\hat{\Phi}(\mathbf{k})$ is found by dividing the density contrast $\hat{\delta}(\mathbf{k})$ by $-k^2$. The potential in real space is then recovered via an inverse FFT. For these operations, we leveraged the cuFFT library through its PyTorch interface (`torch.fft.rfftn` and `torch.fft.irfftn`).

\subsubsection{Time integration}
Particle trajectories are evolved using a standard leapfrog (kick-drift-kick) integrator. The `kick` steps update particle velocities based on the computed gravitational force, while the `drift` step updates their positions based on their velocities.

\subsection{Performance benchmarking}
To quantify the performance of our GPU implementation, we benchmarked it against an equivalent multi-threaded CPU version. The CPU code was written in Python using NumPy for array operations and SciPy (`scipy.fft`) for the FFT-based Poisson solver, configured to use all available CPU cores. We measured the wall-clock time for each major sub-step of the simulation loop (mass assignment, FFT solver, force interpolation) for both implementations. The reported timings are an average over five simulation steps, preceded by a warm-up phase to exclude any one-time kernel compilation overheads.

\subsection{Validation and power spectrum estimation}
We validated the physical fidelity of our pipeline by comparing the matter power spectrum against theoretical predictions. For this test, the initial particle displacements at $z=127$ were linearly evolved to $z=0$ by scaling them with the linear growth factor, $D(z=0)/D(z=127)$, calculated using CAMB.

The matter power spectrum, $P(k)$, was estimated from the resulting $z=0$ particle snapshot using the following procedure. First, a density field was generated by assigning particles to a $512^3$ grid using the CIC method. This field was then converted to the density contrast, $\delta$. We computed the 3D Fourier transform of the density contrast field and calculated the raw power by averaging the squared amplitudes of the Fourier modes, $|\hat{\delta}(\mathbf{k})|^2$, in spherical shells of constant wavenumber, $k = |\mathbf{k}|$.

To obtain the final power spectrum, two corrections were applied. First, we deconvolved the effect of the CIC mass assignment scheme by dividing the raw power by the square of the CIC window function, $W(\mathbf{k})$:
\begin{equation}
    W(\mathbf{k}) = \text{sinc}^2\left(\frac{k_x \Delta x}{2}\right) \text{sinc}^2\left(\frac{k_y \Delta y}{2}\right) \text{sinc}^2\left(\frac{k_z \Delta z}{2}\right)
\end{equation}
where $\Delta x, \Delta y, \Delta z$ are the grid cell sizes. Second, we subtracted the shot noise contribution, which is constant and equal to $1/\bar{n}$, where $\bar{n}$ is the mean particle number density. The resulting $P(k)$ was then compared to the linear theory prediction from CAMB at $z=0$.

\end{document}
                